\documentclass[aspectratio=169,10pt]{beamer}

\usepackage[T1]{fontenc}
\usepackage{lmodern}
\usepackage{amsmath}
\usepackage{booktabs}
\usepackage{graphicx}
\usepackage{xcolor}
\usepackage{array}

\definecolor{cmsblue}{RGB}{0,92,170}
\definecolor{cmsred}{RGB}{220,35,55}
\definecolor{cmsgreen}{RGB}{0,145,85}
\definecolor{cmsamber}{RGB}{185,110,0}
\setbeamercolor{structure}{fg=cmsblue}
\setbeamercolor{frametitle}{fg=black}
\setbeamercolor{title}{fg=black}
\setbeamercolor{alerted text}{fg=cmsred}
\setbeamertemplate{navigation symbols}{}
\setbeamertemplate{itemize item}{\color{cmsblue}$\blacktriangleright$}
\setbeamertemplate{footline}{%
  \leavevmode%
  \hbox{%
    \begin{beamercolorbox}[wd=.82\paperwidth,ht=2.6ex,dp=1.1ex,leftskip=1.0em]{author in head/foot}%
      \color{gray}\scriptsize CMS Work in progress
    \end{beamercolorbox}%
    \begin{beamercolorbox}[wd=.18\paperwidth,ht=2.6ex,dp=1.1ex,rightskip=1.0em plus1fil]{date in head/foot}%
      \color{gray}\scriptsize\insertframenumber/\inserttotalframenumber
    \end{beamercolorbox}}%
  \vskip0pt}
\setbeamertemplate{headline}{%
  \color{cmsblue}\rule{\paperwidth}{1.2pt}}
\setbeamersize{text margin left=0.55cm,text margin right=0.55cm}
\graphicspath{{output/flavorMatrix/}}

\newcommand{\plotbox}[2][0.93\textheight]{%
  \centering\includegraphics[height=#1,keepaspectratio]{#2}}
\newcommand{\cmark}{\textcolor{cmsgreen}{\textbf{+}}}
\newcommand{\xmark}{\textcolor{cmsred}{\textbf{-}}}

\title{Hybrid flavor tagging and response in all-pairs Z+jet}
\subtitle{ParticleNet QvG, parallel barrel response, and a four-flavor matrix fit}
\author{M. Voutilainen}
\date{Run 2024I and Summer24 DY -- one-file local validation}

\begin{document}

\begin{frame}[plain]
  \titlepage
  \vfill
  \begin{center}
    \color{cmsamber}\small Numerical validation only; not a precision flavor correction
  \end{center}
\end{frame}

\begin{frame}{The hybrid tag has four reconstructed categories}
\begin{columns}[T,onlytextwidth]
\column{0.48\textwidth}
\begin{block}{Sequential definition}
\begin{align*}
 b &: \mathrm{UParT\ CvB}<0.5,\\
 c &: \mathrm{CvB}\geq0.5,\quad \mathrm{CvL}\geq0.5,\\
 uds &: \mathrm{CvB}\geq0.5,\quad \mathrm{CvL}<0.5,\quad
        \mathrm{PNet\ QvG}\geq0.3,\\
 g &: \mathrm{CvB}\geq0.5,\quad \mathrm{CvL}<0.5,\quad
        \mathrm{PNet\ QvG}<0.3.
\end{align*}
\end{block}
\column{0.50\textwidth}
\begin{itemize}
\item ParticleNet replaces the time-dependent UParT QvG node.
\item The QvG decision is applied only after the b and c vetoes.
\item The empty reconstructed-undefined category is removed from fits.
\item Truth $d/u+s$ is combined into $uds$; truth undefined is combined with $g$ after checking its discriminator shapes.
\item Four measured tags now constrain four truth groups.
\end{itemize}
\end{columns}
\end{frame}

\begin{frame}{A QvG threshold of 0.3 favors gluon purity}
\begin{columns}[T,onlytextwidth]
\column{0.60\textwidth}
\plotbox[0.78\textheight]{qvg_working_point_scan.pdf}
\column{0.38\textwidth}
\small
\begin{tabular}{lcc}
\toprule
 & 0.3 & 0.5\\
\midrule
$g$ purity & 0.631 & 0.521\\
$g$ efficiency & 0.600 & 0.808\\
$uds$ purity & 0.775 & 0.861\\
$uds$ efficiency & 0.781 & 0.592\\
\bottomrule
\end{tabular}
\vspace{0.5em}
\begin{alertblock}{Working point choice}
The lower threshold trades some gluon efficiency for an 11 percentage-point gain in gluon purity, while recovering substantial uds efficiency.
\end{alertblock}
\end{columns}
\end{frame}

\begin{frame}{The score cube is now separated from its legend}
\plotbox[0.86\textheight]{tagger_cube_allflavors.pdf}
\end{frame}

\begin{frame}{One-dimensional scores show the dominant handles}
\begin{columns}[c,onlytextwidth]
\column{0.33\textwidth}\includegraphics[width=\linewidth]{tagger_CvB_trueflavor.pdf}
\column{0.33\textwidth}\includegraphics[width=\linewidth]{tagger_CvL_trueflavor.pdf}
\column{0.33\textwidth}\includegraphics[width=\linewidth]{tagger_QvG_trueflavor.pdf}
\end{columns}
\end{frame}

\begin{frame}{Pairwise axes and stored discriminator coordinates now agree}
\begin{columns}[c,onlytextwidth]
\column{0.33\textwidth}\includegraphics[width=\linewidth]{tagger2d_CvL_CvB_allflavors.pdf}
\column{0.33\textwidth}\includegraphics[width=\linewidth]{tagger2d_QvG_CvB_allflavors.pdf}
\column{0.33\textwidth}\includegraphics[width=\linewidth]{tagger2d_QvG_CvL_allflavors.pdf}
\end{columns}
\vspace{-0.2em}
\centering\scriptsize Horizontal coordinates are the second stored discriminator; vertical coordinates are the first. Black lines show the actual 0.5/0.3 category boundaries.
\end{frame}

\begin{frame}{Tagging efficiencies: MC, inferred data, and their ratio}
\begin{columns}[c,onlytextwidth]
\column{0.33\textwidth}\includegraphics[width=\linewidth]{efficiency_mc.pdf}
\column{0.33\textwidth}\includegraphics[width=\linewidth]{efficiency_data_inferred.pdf}
\column{0.33\textwidth}\includegraphics[width=\linewidth]{efficiency_transition_sf.pdf}
\end{columns}
\vspace{-0.2em}
\centering\scriptsize Data transitions are the minimum-KL solution with fixed MC truth fractions, not independent truth-labeled measurements.
\end{frame}

\begin{frame}{Purities are the relevant mixture weights for response}
\begin{columns}[c,onlytextwidth]
\column{0.33\textwidth}\includegraphics[width=\linewidth]{purity_mc.pdf}
\column{0.33\textwidth}\includegraphics[width=\linewidth]{purity_data_inferred.pdf}
\column{0.33\textwidth}\includegraphics[width=\linewidth]{purity_ratio_data_over_mc.pdf}
\end{columns}
\vspace{-0.2em}
\centering\scriptsize $P(f\mid t)$ enters each reconstructed-tag response. The b tag is very pure; c and g remain the most composition-sensitive rows.
\end{frame}

\begin{frame}{Transition responses expose sparse heavy-flavor cells}
\begin{columns}[c,onlytextwidth]
\column{0.33\textwidth}\includegraphics[width=\linewidth]{response_transition_mc.pdf}
\column{0.33\textwidth}\includegraphics[width=\linewidth]{response_transition_data_estimated.pdf}
\column{0.33\textwidth}\includegraphics[width=\linewidth]{response_transition_ratio_data_over_mc.pdf}
\end{columns}
\vspace{-0.2em}
\centering\scriptsize The data map assumes one data/MC response factor per truth flavor, common to all reconstructed tags. Low-purity cells can have large MC template fluctuations.
\end{frame}

\begin{frame}{Pure-parallel HDM gives a full-rank four-flavor fit}
\begin{columns}[T,onlytextwidth]
\column{0.51\textwidth}
\begin{align*}
P_{\mathrm{data}}(f\mid t) &=
\frac{W_{\mathrm{data}}(t,f)}{\sum_{f'}W_{\mathrm{data}}(t,f')},\\
R_{\mathrm{data}}(t) &= \sum_f P_{\mathrm{data}}(f\mid t)
R_{\mathrm{MC}}(t,f)\,k_f,\\
R_{\mathrm{HDM}} &= \frac{m_0-m_n-m_u}
{1-m_n/R_n-m_u/R_u}.
\end{align*}
\column{0.47\textwidth}
\begin{block}{Current choices}
\begin{itemize}
\item $|\eta_{jet}|<1.3$ and pure parallel window
\item $p_{T,Z}>30$ GeV
\item $R_n=1.00$, $R_u=0.92$
\item merged component means, never event-wise HDM
\item four reco rows, four truth columns
\item rank $4/4$, condition number 7.37
\end{itemize}
\end{block}
\end{columns}
\end{frame}

\begin{frame}{The local file determines only percent-level flavor residuals}
\begin{columns}[c,onlytextwidth]
\column{0.55\textwidth}
\includegraphics[width=\linewidth]{flavor_response_residual.pdf}
\column{0.43\textwidth}
\small
\begin{tabular}{lcc}
\toprule
Flavor & data/MC & conditional stat.\\
\midrule
$uds$ & 0.986 & 5.1\%\\
$c$   & 0.875 & 10.0\%\\
$b$   & 0.916 & 15.6\%\\
$g$   & 1.072 & 9.5\%\\
\bottomrule
\end{tabular}
\vspace{0.6em}
\begin{alertblock}{Interpretation}
The c and b central values are not yet precision results. The local sample is two orders of magnitude away from a 0.1\% statistical target even before model systematics.
\end{alertblock}
\end{columns}
\end{frame}

\begin{frame}{Only two local pT bins pass the stability gate}
\begin{columns}[c,onlytextwidth]
\column{0.57\textwidth}
\includegraphics[width=\linewidth]{flavor_response_residual_vs_pt.pdf}
\column{0.41\textwidth}
\scriptsize
\begin{tabular}{lrrl}
\toprule
$p_T$ (GeV) & rank & condition & plot\\
\midrule
30--40  & 4 & $5.8\times10^7$ & rejected\\
40--60  & 4 & 7.37 & shown\\
60--85  & 4 & 4.49 & shown\\
85--125 & 3 & 4.69 & rejected\\
125--180& 3 & $2.7\times10^7$ & rejected\\
180--250& 3 & $1.2\times10^7$ & rejected\\
250--400& 2 & 4.88 & rejected\\
\bottomrule
\end{tabular}
\vspace{0.5em}
The TSV output retains every bin. The figure requires rank 4 and condition number below 100.
\end{columns}
\end{frame}

\begin{frame}{Absolute response uncertainties identify the bottleneck}
\begin{columns}[c,onlytextwidth]
\column{0.56\textwidth}
\centering
\includegraphics[width=0.92\linewidth]{response_uncertainty_components.pdf}
\column{0.42\textwidth}
\small
\begin{itemize}
\item Data response statistics dominate this one-file test.
\item MC response-template statistics contribute 1--3\%.
\item A multinomial MC-purity estimate is below 0.4\%.
\item Independent 1\% truth-fraction changes shift the fitted responses by less than 0.05\% in this conditional model.
\end{itemize}
\begin{alertblock}{Not yet a full covariance}
Generator modeling, tagger-shape systematics, correlations, JEC/JER, pileup, and the signal/sideband likelihood remain to be added.
\end{alertblock}
\end{columns}
\end{frame}

\begin{frame}{Semileptonic heavy-flavor decays require a separate truth control}
\begin{columns}[T,onlytextwidth]
\column{0.48\textwidth}
\begin{block}{Physical effect}
Neutrinos from semileptonic c- and b-hadron decays lower the visible DB/MPF-HDM response. Depending on the particle-level jet definition, that loss need not appear in the nominal no-neutrino truth response.
\end{block}
\column{0.50\textwidth}
\begin{block}{Next-run control}
\begin{itemize}
\item identify heavy-hadron descendant neutrinos with GenPart ancestry;
\item store neutrino $p_T$ fraction and $R_{gen+\nu}/R_{gen}$ by truth/reco cell;
\item retain an ambiguous ancestry category;
\item separate decay-neutrino physics from detector response and tagging migrations.
\end{itemize}
\end{block}
\end{columns}
\end{frame}

\begin{frame}{Recommended path before the next full Condor campaign}
\begin{enumerate}
\item Keep PNet QvG at 0.3 and validate the purity scan on the full 2024I sample.
\item Use pure-parallel barrel response as the current central flavor fit; retain transverse profiles as a pileup systematic and forward control.
\item Add heavy-decay-neutrino and direct-versus-splitting truth controls before interpreting c/b residuals physically.
\item Require full rank and a bounded condition number separately in every reported pT bin.
\item Move to the full 2024 sample only after closure with injected flavor-response shifts and varied truth fractions.
\end{enumerate}
\vfill
\begin{center}
\color{cmsgreen}\textbf{The present local sample validates the method and exposes its precision limit; it does not yet measure 0.1\% flavor corrections.}
\end{center}
\end{frame}

\end{document}
